\documentclass[a4paper,oneside,12pt]{report}
% Package yang diperlukan
\usepackage{graphicx}      % Untuk memasukkan gambar
\usepackage{titling}       % Untuk kustomisasi halaman judul
\usepackage{geometry}      % Untuk mengatur margin halaman
\usepackage{listings}      % Untuk menampilkan blok kode
\usepackage{xcolor}        % Untuk menggunakan warna
\usepackage{float}         % Untuk kontrol penempatan float (seperti gambar)
\usepackage{hyperref}      % Untuk membuat hyperlink (URL, referensi)
\usepackage{indentfirst}   % Untuk memberi indentasi pada paragraf pertama setiap seksi
\usepackage{url}           % Untuk format URL
\usepackage{fancyhdr}      % Untuk kustomisasi header dan footer
\usepackage{tocloft}       % Untuk kustomisasi Daftar Isi
\usepackage{bookmark}      % Untuk manajemen bookmark PDF yang lebih baik
\usepackage{amsmath}       % Untuk notasi matematika
\usepackage[most]{tcolorbox} % Untuk box output terminal

% Paket untuk tabel agar sesuai format referensi
\usepackage{array}
\usepackage{booktabs}
\usepackage{tabularx}
\usepackage{multirow}
\usepackage{colortbl}

% Macro untuk header tabel
\newcommand{\thead}[1]{\textcolor{white}{\textbf{#1}}}

%======================================================================
% PENGATURAN GEOMETRI HALAMAN
%======================================================================
\geometry{left=3cm, right=2.5cm, top=2cm, bottom=2.5cm}
\sloppy

%======================================================================
% PENGATURAN WARNA UNTUK BLOK KODE
%======================================================================
\definecolor{codegreen}{rgb}{0,0.6,0}
\definecolor{codegray}{rgb}{0.5,0.5,0.5}
\definecolor{codepurple}{rgb}{0.58,0,0.82}
\definecolor{backcolour}{rgb}{0.95,0.95,0.92}
\definecolor{darkblue}{rgb}{0,0,0.6}
\definecolor{outputbg}{RGB}{241,245,249}
\definecolor{outputborder}{RGB}{203,213,225}

%======================================================================
% PENGATURAN BLOK KODE (LISTINGS)
%======================================================================
\lstdefinestyle{commonstyle}{
    backgroundcolor=\color{backcolour},
    breakatwhitespace=false,
    breaklines=true,
    captionpos=b,
    keepspaces=true,
    numbers=left,
    numbersep=5pt,
    showspaces=false,
    showstringspaces=false,
    showtabs=false,
    tabsize=2,
    frame=single,
    basicstyle=\ttfamily\footnotesize,
    numberstyle=\tiny\color{codegray},
    columns=flexible,
    postbreak=\mbox{\textcolor{red}{$\hookrightarrow$}\space},
    breakindent=0pt,
    breakautoindent=false
}

\lstdefinestyle{pythonstyle}{
    style=commonstyle,
    language=Python,
    commentstyle=\color{codegreen},
    keywordstyle=\color{darkblue},
    stringstyle=\color{codepurple},
    emph={import,from,class,def,for,while,if,else,elif,try,except,finally,as,lambda,return},
    emphstyle=\color{darkblue}
}

% ── Box Format Output / Terminal Box (Verbatim via tcblisting) ───────────
\newtcblisting{outputbox}{
    enhanced,
    colback=outputbg,
    colframe=outputborder,
    boxrule=0.6pt,
    arc=1.2mm,
    left=3mm,
    right=3mm,
    top=1.2mm,
    bottom=1.2mm,
    listing only,
    listing options={
        basicstyle=\ttfamily\footnotesize,
        breaklines=true,
        numbers=none,
        showstringspaces=false
    }
}

% ── Command Placeholder Screenshot ────────────────────────────────────────
\newcommand{\screenshotplaceholder}[2]{%
    \begin{figure}[H]
        \centering
        % \includegraphics[width=0.9\textwidth]{#1}
        \begin{tcolorbox}[
            enhanced,
            colback=gray!5,
            colframe=gray!30,
            arc=1.2mm,
            width=0.9\textwidth,
            height=4.5cm,
            valign=center,
            halign=center
        ]
            \textcolor{gray!60}{\fontsize{10}{12}\selectfont\ttfamily [Screenshot: #2]}
        \end{tcolorbox}
        \caption{#2}
        \label{fig:#1}
    \end{figure}
}

\newcommand{\screenshotimage}[3]{%
    \begin{figure}[H]
        \centering
        \includegraphics[width=0.9\textwidth]{#1}
        \caption{#2}
        \label{fig:#3}
    \end{figure}
}

%======================================================================
% PENGATURAN LAIN-LAIN
%======================================================================
\renewcommand{\figurename}{Gambar}
\renewcommand{\thefigure}{\arabic{figure}}

\hypersetup{
    colorlinks=true,
    linkcolor=black,
    filecolor=black,
    urlcolor=black,
    citecolor=black,
    pdftitle={Laporan Praktikum Penambangan Data - Pertemuan I},
    pdfpagemode=FullScreen,
    breaklinks=true
}

\pagestyle{fancy}
\fancyhf{}
\fancyfoot[C]{\thepage}
\renewcommand{\headrulewidth}{0pt}

\fancypagestyle{plain}{
    \fancyhf{}
    \fancyfoot[C]{\thepage}
    \renewcommand{\headrulewidth}{0pt}
}

%======================================================================
% PENGATURAN STRUKTUR BAB DAN DAFTAR ISI
%======================================================================
\renewcommand{\contentsname}{Daftar Isi}
\renewcommand{\bibname}{Daftar Pustaka}

\renewcommand{\thechapter}{}
\renewcommand{\chaptername}{}

\renewcommand{\thesection}{\arabic{chapter}.\arabic{section}}
\renewcommand{\thesubsection}{\thesection.\arabic{subsection}}

\setcounter{secnumdepth}{3}
\setcounter{tocdepth}{3}

\renewcommand{\cftchapdotsep}{\cftdotsep}
\renewcommand{\cftchapleader}{\cftdotfill{\cftchapdotsep}}
\setlength{\cftchapindent}{0em}
\setlength{\cftsecindent}{2em}
\setlength{\cftsubsecindent}{4em}
\setlength{\cftchapnumwidth}{0em}
\setlength{\cftsecnumwidth}{2.5em}
\setlength{\cftsubsecnumwidth}{3.2em}
\renewcommand{\cftchappresnum}{}
\renewcommand{\cftchapaftersnum}{}

\newcommand{\tocboldsection}[1]{%
  \protect\renewcommand{\protect\cftchapfont}{\protect\bfseries}%
  #1%
  \protect\renewcommand{\protect\cftchapfont}{}%
}

%======================================================================
% AWAL DOKUMEN
%======================================================================
\begin{document}

\newgeometry{left=3cm, right=2.5cm, top=2.5cm, bottom=3cm}
\begin{titlingpage}
\begin{center}
\vspace*{0.5cm}
{\large \textbf{\MakeUppercase{Laporan Praktikum}}} \\[0.4cm]
{\large \textbf{\MakeUppercase{Praktikum Penambangan Data}}} \\[0.4cm]
{\large \textbf{\MakeUppercase{Pertemuan I}}} \\[0.4cm]
{\large \textbf{\MakeUppercase{Pengenalan NumPy dan Pandas}}} \\[1.5cm]

\includegraphics[height=7cm]{lambang ugm.png} \\[1cm]

{\normalsize \textbf{Disusun oleh:}} \\[0.4cm]
\begin{tabular}{@{}l l @{}}
  \textbf{Nama}           & : Januarsyah Akbar \\
  \textbf{NIM}            & : 24/535846/SV/24314 \\
  \textbf{Kelas}          & : B2 \\
  \textbf{Dosen Pengampu} & : Dr. Ganjar Alfian, S.T., M.Eng. \\
                          & \phantom{:} Dr. Imam Fahrurrozi, S.T., M.Cs.
\end{tabular} \\[1.2cm]

{\normalsize \textbf{PROGRAM STUDI D-IV}} \\[0.2cm]
{\normalsize \textbf{TEKNOLOGI REKAYASA PERANGKAT LUNAK}} \\[0.2cm]
{\normalsize \textbf{DEPARTEMEN TEKNIK ELEKTRO DAN INFORMATIKA}} \\[0.2cm]
{\normalsize \textbf{SEKOLAH VOKASI}} \\[0.2cm]
{\normalsize \textbf{UNIVERSITAS GADJAH MADA}} \\[0.2cm]
{\normalsize \textbf{YOGYAKARTA}} \\[0.2cm]
{\normalsize \textbf{2026}} \\[1.2cm]
\end{center}
\end{titlingpage}
\restoregeometry

\thispagestyle{empty}
\clearpage
\stepcounter{page}
\thispagestyle{empty}

% --- Daftar Isi ---
\addcontentsline{toc}{chapter}{Daftar Isi}
\tableofcontents
\clearpage

\pagenumbering{arabic}
\setcounter{page}{3}

%======================================================================
% BAB I: TUJUAN PRAKTIKUM
%======================================================================
\chapter[Tujuan Praktikum]{Tujuan Praktikum}
Tujuan pelaksanaan praktikum pada Pertemuan I adalah sebagai berikut:
\begin{enumerate}
    \item Mahasiswa mampu memahami konsep dasar pustaka \textit{NumPy} serta penggunaan struktur array $n$-dimensi (\textit{ndarray}) dan berbagai operasi manipulasi array.
    \item Mahasiswa mengetahui berbagai teknik manipulasi, eksplorasi, dan transformasi data tabular menggunakan pustaka \textit{Pandas}, termasuk pembuatan dan restrukturisasi \textit{DataFrame}.
    \item Mahasiswa dapat menerapkan integrasi \textit{NumPy} dan \textit{Pandas} untuk mengolah, menyaring, mengagregasi, serta menganalisis dataset sebagai pondasi komputasi pada penambangan data (\textit{data mining}).
\end{enumerate}
\clearpage

%======================================================================
% BAB II: DASAR TEORI
%======================================================================
\chapter{Dasar Teori}

\section{NumPy (Numerical Python)}
\textit{NumPy} adalah pustaka inti dalam ekosistem komputasi ilmiah berbasis Python yang menyediakan objek array multidimensi berkinerja tinggi, dikenal sebagai \texttt{ndarray} (\textit{N-dimensional array}). Dibandingkan dengan tipe data bawaan \textit{list} pada Python murni, \textit{ndarray} menyimpan data dalam blok memori yang kontinu secara homogen (\textit{contiguous homogeneous memory layout}), memungkinkan operasi matematika vektorisasi (\textit{vectorization}) dieksekusi secara cepat tanpa beban komputasi perulangan eksplisit \cite{ref1}.

Fitur-fitur fundamental \textit{NumPy} meliputi:
\begin{itemize}
    \item \textbf{Array Creation \& Attributes}: Pembuatan array satu dimensi hingga multidimensi serta inspeksi bentuk dimensi melalui atribut \texttt{shape}, tipe data elemen \texttt{dtype}, dan ukuran data \texttt{size}.
    \item \textbf{Indexing dan Slicing}: Pengambilan elemen spesifik berdasarkan indeks berbasis nol (\textit{zero-based indexing}) serta pemotongan subset array menggunakan sintaks rentang \texttt{[start:stop:step]} tanpa perlu menyalin data fisik (\textit{memory view}).
    \item \textbf{Reshaping \& Flattening}: Pengubahan geometri array tanpa mengubah nilai elemen melalui metode \texttt{reshape()}, serta perataan ke 1D melalui \texttt{ravel()} atau \texttt{reshape(-1)}.
    \item \textbf{Array Operations}: Operasi aritmatika elemen-ke-elemen (\textit{element-wise operations}), penggabungan (\texttt{concatenate}), pencarian berbasis kondisi (\texttt{np.where}), dan fungsi statistik (\texttt{sum}, \texttt{mean}, \texttt{std}, \texttt{prod}).
\end{itemize}

\section{Pandas}
\textit{Pandas} adalah pustaka analisis dan manipulasi data terstruktur yang dibangun di atas \textit{NumPy} \cite{ref2}. \textit{Pandas} menyediakan dua struktur data utama:
\begin{enumerate}
    \item \textbf{Series}: Struktur data satu dimensi berlabel yang mampu menampung berbagai tipe data (integer, string, float, dsb.).
    \item \textbf{DataFrame}: Struktur data dua dimensi berbentuk tabel dua sumbu (\textit{rows} dan \textit{columns}) dengan indeks fleksibel.
\end{enumerate}

Operasi krusial pada \textit{Pandas} meliputi:
\begin{itemize}
    \item \textbf{Data Cleaning \& Transformation}: Penambahan atau penghapusan kolom (\texttt{drop}), pengubahan tipe data, penanganan nilai hilang (\textit{missing values}), dan pemfilteran berbasis kondisi boolean.
    \item \textbf{Grouping \& Aggregation}: Operasi \textit{split-apply-combine} menggunakan metode \texttt{groupby()} yang dikombinasikan dengan \texttt{agg()} untuk menghitung berbagai metrik ringkasan (misalnya \textit{mean}, \textit{max}, \textit{min}, \textit{std}, \textit{count}) pada kelompok data tertentu.
    \item \textbf{Relational Operations (Merge, Join, Concat)}: Pengbubungan beberapa tabel berdasarkan kolom kunci tertentu (\texttt{merge}) atau penggabungan berorientasi sumbu horizontal/vertikal (\texttt{concat}).
    \item \textbf{Reshaping}: Transformasi bentuk tabel melalui pembuatan tabel silang (\texttt{pivot\_table}) dan restrukturisasi dari format lebar (\textit{wide}) ke panjang (\textit{long}) menggunakan metode \texttt{melt}.
    \item \textbf{Functional Application}: Penerapan ekspresi logika kustom pada seluruh baris atau kolom melalui fungsi \texttt{apply()} dengan ekspresi \texttt{lambda}.
\end{itemize}
\clearpage

%======================================================================
% BAB III: HASIL DAN PEMBAHASAN
%======================================================================
\chapter{Hasil dan Pembahasan}

\section{Langkah Praktikum}

\subsection{Langkah Percobaan NumPy}
Di bawah ini merupakan penjelasan detail mengenai langkah-langkah praktikum dasar penggunaan pustaka \textit{NumPy}:

\begin{enumerate}
    \item \textbf{Pembuatan Array 1D, 2D, dan 3D beserta Pengecekan .shape} \\
    \textit{Deskripsi:} Menginisialisasi array satu dimensi (1D), dua dimensi (2D), dan tiga dimensi (3D) menggunakan \texttt{np.array()}, lalu memeriksa dimensi geometri array menggunakan atribut \texttt{.shape}. \\
    \textit{Input Kode:}
\begin{lstlisting}[style=pythonstyle, caption=Inisialisasi dan Cek Shape NumPy]
import numpy as np

# Array 1 Dimensi
arr1d = np.array([1, 2, 3, 4, 5])
print("Array 1D:", arr1d)
print("Shape 1D:", arr1d.shape)
print("Tipe Data:", type(arr1d))
print("-" * 30)

# Array 2 Dimensi (2 baris, 3 kolom)
arr2d = np.array([[1, 2, 3], [4, 5, 6]])
print("Array 2D:")
print(arr2d)
print("Shape 2D:", arr2d.shape)
print("-" * 30)

# Array 3 Dimensi (3 blok, 2 baris, 4 kolom)
arr3d = np.array([
    [[1, 2, 3, 4], [5, 6, 7, 8]],
    [[9, 10, 11, 12], [13, 14, 15, 16]],
    [[17, 18, 19, 20], [21, 22, 23, 24]]
])
print("Array 3D:")
print(arr3d)
print("Shape 3D:", arr3d.shape)
\end{lstlisting}
    \textit{Penjelasan Kode:} Kode di atas mendefinisikan array NumPy dari \textit{list} bersarang. Parameter array diatur sesuai dengan orientasi sumbu baris dan kolom. Atribut \texttt{shape} digunakan untuk memperoleh dimensi dari \textit{ndarray}. \\
    \textit{Hasil Output:} \\
    \screenshotimage{"gambar/1.1 Pembuatan Array 1D, 2D, dan 3D beserta Pengecekan .shape".png}{Hasil Eksekusi Pembuatan Array dan Pengecekan Shape}{ss-numpy-shape}
    \textit{Penjelasan Output:} Array 1D menghasilkan struktur vektor baris berisi 5 elemen dengan bentuk \texttt{(5,)}. Array 2D menghasilkan bentuk \texttt{(2, 3)} (2 baris dan 3 kolom). Array 3D membentuk dimensi matriks berlapis dengan konfigurasi \texttt{(3, 2, 4)} yang berarti terdapat 3 blok matriks berukuran $2\times4$.

    \item \textbf{Indexing dan Slicing Array Multidimensi} \\
    \textit{Deskripsi:} Mengakses elemen array pada koordinat indeks tertentu serta memotong subset array menggunakan notasi titik dua \texttt{[start:stop:step]}. \\
    \textit{Input Kode:}
\begin{lstlisting}[style=pythonstyle, caption=Indexing dan Slicing NumPy]
# Indexing 1D & 2D
print("Elemen ke-0 arr1d:", arr1d[0])
print("Elemen baris 0 kolom 1 arr2d:", arr2d[0, 1])
print("Elemen blok 1 baris 1 kolom 2 arr3d:", arr3d[1, 1, 2])

# Slicing Array
print("\nSlicing arr1d (indeks 1 sampai 3):", arr1d[1:4])
print("Slicing arr2d (semua baris, kolom 1-2):")
print(arr2d[:, 1:])
\end{lstlisting}
    \textit{Penjelasan Kode:} Koordinat indexing menggunakan pemisah koma untuk sumbu baris dan kolom. Pada slicing, tanda \texttt{:} tanpa angka berarti mengambil seluruh elemen pada sumbu tersebut, sedangkan \texttt{1:4} mengambil elemen dari indeks 1 hingga sebelum indeks 4. \\
    \textit{Hasil Output:} \\
    \screenshotimage{"gambar/1.2 Indexing dan Slicing Array Multidimensi".png}{Hasil Eksekusi Indexing dan Slicing NumPy}{ss-numpy-indexing-slicing}
    \textit{Penjelasan Output:} Pengambilan indeks tunggal mengembalikan nilai skalar numerik. Hasil slicing \texttt{arr1d[1:4]} menghasilkan array 1D \texttt{[2, 3, 4]}. Slicing kolom pada \texttt{arr2d} menghasilkan matriks bagian $2\times2$ yang terambil dari kolom kedua dan ketiga.

    \item \textbf{Pembuatan Array dengan Helper Functions} \\
    \textit{Deskripsi:} Menginisialisasi array dengan fungsi instan bawaan NumPy tanpa mendefinisikan nilai elemen secara manual. \\
    \textit{Input Kode:}
\begin{lstlisting}[style=pythonstyle, caption=Pembuatan Array Instan NumPy]
print("Zeros (2x3):")
print(np.zeros((2, 3)))
print("\nOnes (2x2):")
print(np.ones((2, 2)))
print("\nArange (0 sampai 10 kelipatan 2):", np.arange(0, 10, 2))
print("\nLinspace (0 sampai 1 dibagi 5 titik):", np.linspace(0, 1, 5))
print("\nEye / Matriks Identitas (3x3):")
print(np.eye(3))
\end{lstlisting}
    \textit{Penjelasan Kode:} Fungsi \texttt{np.zeros} dan \texttt{np.ones} menerima argumen berbentuk tuple dimensi. Fungsi \texttt{np.arange(start, stop, step)} menghasilkan barisan deret, \texttt{np.linspace(start, stop, num)} membagi rentang menjadi interval yang setara, sedangkan \texttt{np.eye(n)} membuat matriks identitas $n \times n$. \\
    \textit{Hasil Output:} \\
    \screenshotimage{"gambar/1.3 Pembuatan Array dengan Helper Functions (zeros, ones, arange, linspace, eye)".png}{Hasil Eksekusi Helper Functions NumPy}{ss-numpy-helpers}
    \textit{Penjelasan Output:} Fungsi pembantu mempermudah alokasi memori array secara cepat. Nilai desimal standar (\textit{floating-point}) otomatis diimplementasikan untuk \texttt{zeros}, \texttt{ones}, \texttt{linspace}, dan \texttt{eye}.

    \item \textbf{Reshape dan Konversi Tipe Data} \\
    \textit{Deskripsi:} Mengubah dimensi data array menggunakan fungsi \texttt{.reshape()} dan mengonversi tipe data elemen array menggunakan \texttt{.astype()}. \\
    \textit{Input Kode:}
\begin{lstlisting}[style=pythonstyle, caption=Reshape dan Astype NumPy]
data = np.arange(1, 13)
print("Data awal (1D):", data)

# Reshape ke 3x4
reshaped_3x4 = data.reshape(3, 4)
print("\nReshape ke (3, 4):")
print(reshaped_3x4)

# Reshape kembali ke 1D (-1 atau flatten)
flattened = reshaped_3x4.reshape(-1)
print("\nKembali ke 1D:", flattened)

# Konversi Tipe Data
arr_float = data.astype(float)
print("\nKonversi ke Float:", arr_float)
print("Tipe data element:", arr_float.dtype)
\end{lstlisting}
    \textit{Penjelasan Kode:} Pengubahan dimensi dengan \texttt{reshape(r, c)} memerlukan kesesuaian perkalian perkalian baris dan kolom dengan jumlah elemen mula-mula. Sintaks \texttt{reshape(-1)} memberi instruksi otomatis kepada NumPy untuk meratakan bentuk dimensi kembali ke 1D. \\
    \textit{Hasil Output:} \\
    \screenshotimage{"gambar/1.4 Reshape dan Konversi Tipe Data".png}{Hasil Eksekusi Reshape dan Konversi Tipe Data}{ss-numpy-reshape}
    \textit{Penjelasan Output:} Array 1D yang berisi 12 elemen bertransformasi secara mulus menjadi matriks $3\times4$. Setelah dilakukan konversi tipe data \texttt{astype(float)}, data numerik berubah menjadi desimal bertipe \texttt{float64} dengan representasi desimal bertanda titik (\texttt{.}).

    \item \textbf{Operasi Array (Join, Search, Sort, dan Matematika/Statistik)} \\
    \textit{Deskripsi:} Menyatukan beberapa array, melakukan pencarian indeks berdasarkan kondisi tertentu, mengurutkan data, serta menghitung metrik statistik dasar. \\
    \textit{Input Kode:}
\begin{lstlisting}[style=pythonstyle, caption=Operasi Array Lanjutan NumPy]
# Concatenate / Join
a = np.array([1, 2, 3])
b = np.array([4, 5, 6])
print("Join 1D:", np.concatenate((a, b)))

# NumPy Random
np.random.seed(42)
random_arr = np.random.randint(0, 100, size=(3, 3))
print("\nArray Random (3x3):")
print(random_arr)

# Search dengan np.where
indices = np.where(random_arr > 50)
print("\nIndeks elemen > 50:", indices)
print("Nilai elemen > 50:", random_arr[indices])

# Sorting
print("\nSorted array (per baris):")
print(np.sort(random_arr, axis=1))

# Operasi Statistik / Matematika
print("\nTotal Sum:", np.sum(random_arr))
print("Sum per Kolom (axis 0):", np.sum(random_arr, axis=0))
print("Mean / Rata-rata:", np.mean(random_arr))
\end{lstlisting}
    \textit{Penjelasan Kode:} Pembangkitan data menggunakan \texttt{np.random.randint()} dengan penentuan seed deterministik. Pencarian bersyarat memanfaatkan \texttt{np.where()} yang menghasilkan tuple array indeks koordinat. Operasi statistik dilakukan secara terarah pada baris atau kolom menggunakan parameter \texttt{axis}. \\
    \textit{Hasil Output:} \\
    \screenshotimage{"gambar/1.5 Operasi Array (Join, Split, Search, Sort, Random, dan Matematika)".png}{Hasil Eksekusi Operasi Array Lanjutan}{ss-numpy-operations}
    \textit{Penjelasan Output:} Fungsi \texttt{np.where} mengembalikan elemen bernilai \texttt{[51, 92, 71, 60, 82, 86, 74]}. Pengurutan dengan parameter \texttt{axis=1} mengatur ulang elemen baris agar berurutan secara menaik. Perhitungan akumulasi penjumlahan menghasilkan nilai total \texttt{550}, penjumlahan per kolom (\texttt{axis=0}) menghasilkan \texttt{[204, 238, 108]}, dan nilai rata-rata keseluruhan didapat sebesar \texttt{61.11}.
\end{enumerate}

\subsection{Langkah Percobaan Pandas}
Di bawah ini merupakan penjelasan detail mengenai langkah-langkah praktikum dasar penggunaan pustaka \textit{Pandas}:

\begin{enumerate}
    \item \textbf{Pembuatan DataFrame, Penambahan, dan Penghapusan Kolom} \\
    \textit{Deskripsi:} Membuat objek \textit{DataFrame} dari data berformat \textit{dictionary}, menambahkan atribut kolom baru berbasis operasi aritmatika, serta membuang kolom yang tidak diperlukan. \\
    \textit{Input Kode:}
\begin{lstlisting}[style=pythonstyle, caption=Pembuatan dan Manipulasi Kolom DataFrame]
# Membuat DataFrame awal
data_mhs = {
    'NIM': ['22/001', '22/002', '22/003', '22/004'],
    'Nama': ['Budi', 'Siti', 'Joko', 'Dewi'],
    'Jurusan': ['TRPL', 'TRPL', 'TRK', 'TRPL'],
    'Nilai_UAS': [85, 92, 78, 88]
}
df_mhs = pd.DataFrame(data_mhs)
print("DataFrame Awal:")
print(df_mhs.to_string(index=False))

# Menambahkan Kolom Baru
df_mhs['Status'] = ['Lulus', 'Lulus', 'Lulus', 'Lulus']
df_mhs['Nilai_Akhir'] = df_mhs['Nilai_UAS'] + 5
print("\nSetelah Menambah Kolom:")
print(df_mhs.to_string(index=False))

# Menghapus Kolom
df_clean = df_mhs.drop(columns=['Status'])
print("\nSetelah Menghapus Kolom Status:")
print(df_clean.to_string(index=False))
\end{lstlisting}
    \textit{Penjelasan Kode:} Inisialisasi awal mengubah struktur dictionary menjadi bentuk tabular. Penambahan kolom langsung memanfaatkan inisialisasi indeks kolom baru. Untuk menghapus kolom secara permanen tanpa mengubah objek orisinal (kecuali menggunakan parameter \texttt{inplace=True}), digunakan metode \texttt{.drop(columns=[...])}. \\
    \textit{Hasil Output:} \\
    \screenshotimage{"gambar/2.1 Pembuatan DataFrame, Penambahan & Penghapusan Kolom".png}{Hasil Eksekusi Manipulasi Kolom DataFrame}{ss-pandas-columns}
    \textit{Penjelasan Output:} Proses inisialisasi berhasil membuat DataFrame berisi 4 baris data. Kolom \texttt{Nilai\_Akhir} terhitung secara dinamis dari penambahan kolom \texttt{Nilai\_UAS} dengan nilai 5. Penghapusan kolom \texttt{Status} sukses dilakukan pada DataFrame \texttt{df\_clean}.

    \item \textbf{Sorting dan Filtering Data} \\
    \textit{Deskripsi:} Menyaring baris data berdasarkan kriteria logika tertentu serta mengurutkan baris data menurut nilai kolom numerik. \\
    \textit{Input Kode:}
\begin{lstlisting}[style=pythonstyle, caption=Filtering dan Sorting Pandas]
# Filtering: Nilai_UAS >= 85 dan Jurusan TRPL
filtered_df = df_mhs[(df_mhs['Nilai_UAS'] >= 85) & (df_mhs['Jurusan'] == 'TRPL')]
print("Hasil Filtering:")
print(filtered_df.to_string(index=False))

# Sorting berdasarkan Nilai_UAS secara Descending
sorted_df = df_mhs.sort_values(by='Nilai_UAS', ascending=False)
print("\nHasil Sorting:")
print(sorted_df.to_string(index=False))
\end{lstlisting}
    \textit{Penjelasan Kode:} Filtering menggunakan gabungan ekspresi boolean \texttt{df[kondisi\_1 \& kondisi\_2]}. Tanda operator bitwise \texttt{\&} wajib digunakan sebagai representasi logika \textit{AND}. Fungsi \texttt{.sort\_values()} menggunakan parameter kolom target dan mengaktifkan pengurutan menurun melalui parameter \texttt{ascending=False}. \\
    \textit{Hasil Output:} \\
    \screenshotimage{"gambar/2.2 Sorting dan Filtering Data".png}{Hasil Eksekusi Filtering dan Sorting Pandas}{ss-pandas-filter-sort}
    \textit{Penjelasan Output:} Hasil filtering menyaring baris sehingga hanya menyisakan data mahasiswa TRPL yang memiliki Nilai UAS minimal 85 (Budi, Siti, dan Dewi). Hasil pengurutan menempatkan data mahasiswa Siti pada urutan teratas dan Joko pada urutan terbawah.

    \item \textbf{Grouping dan Aggregation (\texttt{groupby} \& \texttt{agg})} \\
    \textit{Deskripsi:} Mengelompokkan baris berdasarkan kesamaan kategori pekerjaan serta menghitung beberapa fungsi ringkasan metrik statistik secara kolektif. \\
    \textit{Input Kode:}
\begin{lstlisting}[style=pythonstyle, caption=Pengelompokan dan Agregasi Pandas]
# Data Karyawan
data_karyawan = {
    'Nama': ['Budi', 'Maya', 'Joko', 'Andi', 'Dewi', 'Dian'],
    'Gender': ['Pria', 'Perempuan', 'Pria', 'Pria', 'Perempuan', 'Perempuan'],
    'Pekerjaan': ['Programmer', 'Programmer', 'Data Engineer', 'Manager', 'Manager', 'Manager'],
    'Gaji': [5000, 5200, 4500, 6500, 8000, 6200],
    'Pengalaman_Th': [2, 3, 1, 5, 7, 4]
}
df_emp = pd.DataFrame(data_karyawan)

# Groupby Pekerjaan dengan multiple agregasi
agg_result = df_emp.groupby('Pekerjaan').agg(
    Jumlah_Pegawai=('Nama', 'count'),
    Rata_Gaji=('Gaji', 'mean'),
    Total_Gaji=('Gaji', 'sum'),
    Max_Pengalaman=('Pengalaman_Th', 'max')
).reset_index()

print("Hasil Groupby & Aggregation:")
print(agg_result.to_string(index=False))
\end{lstlisting}
    \textit{Penjelasan Kode:} Pemanggilan \texttt{groupby('Pekerjaan')} memecah dataset, kemudian metode \texttt{.agg()} melakukan ringkasan kalkulasi kustom. Argumen dengan format \texttt{Nama\_Kolom=(Kolom\_Target, Fungsi\_Agregasi)} memperjelas label penamaan atribut baru secara otomatis. \\
    \textit{Hasil Output:} \\
    \screenshotimage{"gambar/2.3 Grouping dan Aggregation (groupby & agg)".png}{Hasil Eksekusi Grouping dan Agregasi}{ss-pandas-groupby-agg}
    \textit{Penjelasan Output:} Ringkasan data menyimpulkan informasi penting per kategori pekerjaan: \textit{Data Engineer} memiliki 1 staf dengan gaji 4500; \textit{Manager} memiliki 3 staf dengan total gaji 20700 dan lama kerja maksimal 7 tahun; \textit{Programmer} beranggotakan 2 staf dengan rata-rata kompensasi sebesar 5100.

    \item \textbf{Penggabungan Data: \texttt{merge}, \texttt{join}, dan \texttt{concat}} \\
    \textit{Deskripsi:} Mengintegrasikan dua entitas tabel yang terpisah menggunakan relasi kunci referensi, serta menumpuk baris data tambahan secara vertikal. \\
    \textit{Input Kode:}
\begin{lstlisting}[style=pythonstyle, caption=Penggabungan DataFrame Relasional]
# DataFrame 1 (Tabel Karyawan)
df_karyawan = pd.DataFrame({
    'Emp_id': ['E01', 'E02', 'E03'],
    'Nama': ['Budi', 'Joko', 'Maya'],
    'Id_Dept': [1, 2, 1]
})

# DataFrame 2 (Tabel Departemen)
df_dept = pd.DataFrame({
    'Id_Dept': [1, 2, 3],
    'Nama_Dept': ['IT Solution', 'Data Analytics', 'Finance']
})

# Merge (Left Join)
merged_df = pd.merge(df_karyawan, df_dept, on='Id_Dept', how='left')
print("Hasil Merge (Left Join):")
print(merged_df.to_string(index=False))

# Concat baris (axis=0)
df_baru = pd.DataFrame({
    'Emp_id': ['E04'],
    'Nama': ['Rina'],
    'Id_Dept': [3]
})
all_karyawan = pd.concat([df_karyawan, df_baru], ignore_index=True)
print("\nHasil Concat Baris:")
print(all_karyawan.to_string(index=False))
\end{lstlisting}
    \textit{Penjelasan Kode:} Metode \texttt{pd.merge} secara default mendeteksi kesamaan kolom atribut via parameter \texttt{on}. Opsi \texttt{how='left'} memastikan seluruh baris pada tabel kiri tetap utuh. Fungsi \texttt{pd.concat} menggabungkan list DataFrame dengan penyesuaian indeks melalui \texttt{ignore\_index=True}. \\
    \textit{Hasil Output:} \\
    \screenshotimage{"gambar/2.4 Penggabungan Data merge, join, dan concat".png}{Hasil Eksekusi Penggabungan Merge dan Concat}{ss-pandas-merging}
    \textit{Penjelasan Output:} Hasil penyatuan relasional menampilkan detail \texttt{Nama\_Dept} yang sesuai dengan relasi nilai \texttt{Id\_Dept}. Pada keluaran \texttt{pd.concat}, baris baru karyawan E04 (Rina) berhasil disisipkan secara vertikal ke baris terbawah dengan penomoran indeks kontinu yang berurutan.

    \item \textbf{Reshaping Data: \texttt{pivot\_table}, \texttt{melt}, dan \texttt{apply(lambda)}} \\
    \textit{Deskripsi:} Melakukan restrukturisasi bentuk tabel dari format silang, format panjang (\textit{unpivot}), serta mengaplikasikan transformasi nilai menggunakan fungsi lambda. \\
    \textit{Input Kode:}
\begin{lstlisting}[style=pythonstyle, caption=Transformasi Reshaping dan Lambda]
# Pivot Table (Rata-rata Gaji per Pekerjaan & Gender)
pivot = pd.pivot_table(df_emp, values='Gaji', index='Pekerjaan', columns='Gender', aggfunc='mean')
print("Pivot Table:")
print(pivot)

# Melt (Wide to Long)
scores = pd.DataFrame({
    'Nama': ['Ali', 'Budi'],
    'Matematika': [80, 90],
    'Fisika': [85, 75]
})
melted = pd.melt(scores, id_vars=['Nama'], value_vars=['Matematika', 'Fisika'], var_name='Mata_Pelajaran', value_name='Nilai')
print("\nHasil Melt (Format Panjang):")
print(melted.to_string(index=False))

# Lambda Function dengan Apply
df_emp['Kategori_Gaji'] = df_emp['Gaji'].apply(lambda x: 'Tinggi' if x >= 6000 else ('Sedang' if x >= 5000 else 'Standar'))
print("\nSetelah Apply Lambda Kategori Gaji:")
print(df_emp[['Nama', 'Pekerjaan', 'Gaji', 'Kategori_Gaji']].to_string(index=False))
\end{lstlisting}
    \textit{Penjelasan Kode:} \texttt{pd.pivot\_table} memproyeksikan kategori kolom \texttt{Gender} menjadi nama kolom baru dengan agregasi nilai gaji rata-rata. Fungsi \texttt{pd.melt} merubah kolom mata pelajaran menjadi representasi nilai baris. Metode \texttt{.apply()} mengevaluasi setiap elemen kolom gaji menggunakan ekspresi inline lambda secara cepat. \\
    \textit{Hasil Output:} \\
    \screenshotimage{"gambar/2.5 Reshaping Data pivot_table, melt, dan apply(lambda)".png}{Hasil Eksekusi Reshaping dan Fungsi Lambda}{ss-pandas-reshaping}
    \textit{Penjelasan Output:} Pivot table menampilkan rerata kompensasi secara bersilang (misal kategori \textit{Data Engineer Perempuan} bernilai \texttt{NaN} karena data kosong). Hasil melting berhasil mereduksi struktur kolom menjadi baris variabel \texttt{Mata\_Pelajaran} dan \texttt{Nilai}. Transformasi lambda berhasil melabeli kategori gaji karyawan sesuai ambang batas yang ditentukan.
\end{enumerate}
\clearpage

\section{Tugas}

\begin{enumerate}
    \item \textbf{Tugas 1: NumPy Reshape dan Boolean Filtering}
    \begin{enumerate}
        \item[1.] \textbf{Deskripsi Tugas.} \\
        Buatlah array NumPy 2D berukuran $4 \times 4$ berisi angka acak (rentang nilai 0 hingga 50), kemudian ubah bentuk dimensinya (\textit{reshape}) menjadi array $2 \times 8$, lalu tampilkan hanya elemen-elemen yang bernilai lebih dari 25.

        \item[2.] \textbf{Implementasi Kode.} \\
        Berikut adalah penulisan blok kode program untuk menyelesaikan instruksi tugas 1:
\begin{lstlisting}[style=pythonstyle, caption=Solusi Tugas 1 - NumPy Reshape dan Filtering]
import numpy as np

# 1. Membuat array 2D 4x4 berisi integer acak 0-50
np.random.seed(42) # Seed untuk reproduktibilitas
arr_soal1 = np.random.randint(0, 51, size=(4, 4))
print("1. Array 2D Asal (4x4):")
print(arr_soal1)
print("   Shape:", arr_soal1.shape)
print("-" * 40)

# 2. Reshape array menjadi 2x8
arr_reshaped = arr_soal1.reshape(2, 8)
print("2. Hasil Reshape (2x8):")
print(arr_reshaped)
print("   Shape:", arr_reshaped.shape)
print("-" * 40)

# 3. Filter dan tampilkan hanya elemen yang > 25
elemen_gt_25 = arr_reshaped[arr_reshaped > 25]
print("3. Elemen-elemen yang bernilai > 25:")
print(elemen_gt_25)
print("   Banyaknya elemen yang memenuhi:", len(elemen_gt_25))
\end{lstlisting}

        \item[3.] \textbf{Hasil Eksekusi.} \\
        \screenshotimage{"gambar/Soal 1 NumPy Reshape & Filtering".png}{Screenshot Hasil Eksekusi Tugas 1}{ss-tugas1}

        \item[4.] \textbf{Analisis dan Pembahasan.} \\
        Pembangkitan matriks dua dimensi berukuran $4 \times 4$ dilakukan dengan memanfaatkan fungsi \texttt{np.random.randint(0, 51, size=(4, 4))} yang menghasilkan bilangan bulat acak pada rentang inklusif $[0, 50]$. Untuk menjamin sifat deterministik sehingga hasil komputasi dapat direplikasi secara konsisten, diatur \texttt{np.random.seed(42)} sebelum proses pengacakan dilakukan. Selanjutnya, transformasi bentuk array dari dimensi $4 \times 4$ menjadi $2 \times 8$ dilakukan melalui fungsi \texttt{reshape(2, 8)}. Operasi ini sangat efisien dengan tingkat kompleksitas waktu $O(1)$ karena tidak memodifikasi data fisik di memori, melainkan hanya mengubah metadata \textit{stride} dan bentuk dimensi internal array. Terakhir, penyaringan elemen bernilai lebih besar dari 25 dilakukan menggunakan mekanisme \textit{boolean masking} (\texttt{arr\_reshaped > 25}). Ekspresi logika tersebut memicu evaluasi tiap elemen menghasilkan array boolean yang memproyeksikan elemen-elemen yang bernilai \texttt{True} menjadi vektor satu dimensi baru dengan total 6 elemen yang lolos penyaringan.
    \end{enumerate}

    \item \textbf{Tugas 2: Pandas Groupby dan Multiple Aggregations}
    \begin{enumerate}
        \item[1.] \textbf{Deskripsi Tugas.} \\
        Buatlah \textit{DataFrame} untuk data karyawan yang berisi kolom \texttt{Nama}, \texttt{Pekerjaan}, dan \texttt{Gaji}. Lalu lakukan operasi \texttt{groupby()} pada kolom \texttt{Pekerjaan} dan hitung nilai \textbf{maksimum}, \textbf{minimum}, serta \textbf{standar deviasi} untuk kolom \texttt{Gaji} menggunakan metode \texttt{agg()}.

        \item[2.] \textbf{Implementasi Kode.} \\
        Berikut adalah penulisan blok kode program untuk menyelesaikan instruksi tugas 2:
\begin{lstlisting}[style=pythonstyle, caption=Solusi Tugas 2 - Pandas Groupby dan Aggregation]
import pandas as pd

# 1. Membuat DataFrame data karyawan
data_soal2 = {
    'Nama': ['Andi', 'Budi', 'Citra', 'Dewi', 'Eko', 'Fani', 'Gita', 'Hadi'],
    'Pekerjaan': ['Programmer', 'Programmer', 'Data Analyst', 'Data Analyst', 'Manager', 'Manager', 'Programmer', 'Data Analyst'],
    'Gaji': [6000000, 6500000, 5500000, 5800000, 9500000, 10000000, 7000000, 6200000]
}

df_karyawan_soal2 = pd.DataFrame(data_soal2)
print("Data Karyawan:")
print(df_karyawan_soal2.to_string())
print("-" * 50)

# 2. Operasi groupby() dan agg() untuk Max, Min, dan Std Deviasi kolom Gaji
hasil_agg = df_karyawan_soal2.groupby('Pekerjaan')['Gaji'].agg(
    Gaji_Maksimum='max',
    Gaji_Minimum='min',
    Standar_Deviasi='std'
).reset_index()

print("\nHasil Agregasi Gaji per Pekerjaan:")
print(hasil_agg.to_string(index=False))
\end{lstlisting}

        \item[3.] \textbf{Hasil Eksekusi.} \\
        \screenshotimage{"gambar/Soal 2 Pandas Groupby & Multiple Aggregations".png}{Screenshot Hasil Eksekusi Tugas 2}{ss-tugas2}

        \item[4.] \textbf{Analisis dan Pembahasan.} \\
        Proses pengolahan data karyawan diawali dengan melakukan pengelompokan baris data menggunakan metode \texttt{groupby('Pekerjaan')} untuk mempartisi dataset menjadi kelompok-kelompok diskrit berdasarkan nama jabatan karyawan, yaitu \textit{Data Analyst}, \textit{Manager}, dan \textit{Programmer}. Di atas data terkelompok tersebut, diaplikasikan fungsi \texttt{.agg()} untuk menghitung beberapa ringkasan metrik statistik secara terpadu dalam satu lintasan komputasi, mencakup gaji maksimum (\texttt{max}), minimum (\texttt{min}), dan standar deviasi (\texttt{std}). Berdasarkan luaran pengolahan, nilai maksimum dan minimum memperlihatkan rentang batas kompensasi pada masing-masing posisi, sedangkan nilai standar deviasi menggambarkan variasi sebaran kompensasi tersebut. Pada pengelompokan ini, jabatan \textit{Programmer} terbukti memiliki sebaran gaji paling variatif dengan deviasi sebesar Rp500.000,00 dibandingkan divisi lainnya. Setelah proses agregasi selesai, dipanggil metode \texttt{reset\_index()} untuk menormalkan kembali kolom \texttt{Pekerjaan} dari semula berkedudukan sebagai indeks baris berjenjang menjadi atribut kolom reguler pada tabel dua dimensi.
    \end{enumerate}

    \item \textbf{Tugas 3: Pandas Relational Merge (Left Join dengan Kolom Kunci Berbeda dan Suffixes)}
    \begin{enumerate}
        \item[1.] \textbf{Deskripsi Tugas.} \\
        Buat dua \textit{DataFrame}: DataFrame karyawan dengan kolom \texttt{Emp\_id}, \texttt{Nama}, dan \texttt{Divisi}, serta DataFrame gaji dengan kolom \texttt{ID}, \texttt{Gaji}, dan \texttt{Status} (masing-masing minimal tiga baris). Setelah itu, gabungkan kedua \textit{DataFrame} tersebut menggunakan \texttt{merge()} dengan parameter \texttt{left\_on='Emp\_id'} dan \texttt{right\_on='ID'}, terapkan metode \textbf{left join}, dan gunakan parameter \texttt{suffixes} untuk membedakan nama kolom yang bertabrakan.

        \item[2.] \textbf{Implementasi Kode.} \\
        Berikut adalah penulisan blok kode program untuk menyelesaikan instruksi tugas 3:
\begin{lstlisting}[style=pythonstyle, caption=Solusi Tugas 3 - Pandas Relational Merge]
import pandas as pd

# 1. Membuat DataFrame Karyawan (Tabel Kiri)
df_info_karyawan = pd.DataFrame({
    'Emp_id': ['EMP01', 'EMP02', 'EMP03', 'EMP04'],
    'Nama': ['Rian', 'Sari', 'Doni', 'Lestari'],
    'Divisi': ['Engineering', 'Data Science', 'Product', 'Marketing'],
    'Status': ['Aktif', 'Aktif', 'Cuti', 'Aktif']
})

# 2. Membuat DataFrame Gaji (Tabel Kanan)
df_gaji_karyawan = pd.DataFrame({
    'ID': ['EMP01', 'EMP02', 'EMP03', 'EMP05'],
    'Gaji': [7500000, 8200000, 6800000, 9000000],
    'Status': ['Tetap', 'Kontrak', 'Tetap', 'Kontrak']
})

print("DataFrame Karyawan (Kiri):")
print(df_info_karyawan.to_string(index=False))
print("\nDataFrame Gaji (Kanan):")
print(df_gaji_karyawan.to_string(index=False))
print("-" * 60)

# 3. Menggabungkan kedua DataFrame menggunakan merge (Left Join, left_on, right_on, suffixes)
df_merged = pd.merge(
    df_info_karyawan,
    df_gaji_karyawan,
    left_on='Emp_id',
    right_on='ID',
    how='left',
    suffixes=('_karyawan', '_kepegawaian')
)

print("\nHasil Merge (Left Join):")
print(df_merged.to_string(index=False))
\end{lstlisting}

        \item[3.] \textbf{Hasil Eksekusi.} \\
        \screenshotimage{"gambar/Soal 3 Pandas Merge (Left Join dengan Key Berbeda & Suffixes)".png}{Screenshot Hasil Eksekusi Tugas 3}{ss-tugas3}

        \item[4.] \textbf{Analisis dan Pembahasan.} \\
        Penggabungan dua dataset dilakukan melalui metode \texttt{pd.merge()} dengan menerapkan mekanisme \textit{left join} (\texttt{how='left'}). Metode ini memastikan seluruh informasi dari DataFrame karyawan di sisi kiri tetap dipertahankan secara utuh, sekalipun baris tersebut tidak memiliki padanan nilai kunci di tabel kanan. Hal ini terlihat pada entitas karyawan dengan ID \texttt{EMP04} (Lestari) yang tetap muncul pada hasil merge, di mana nilai pada kolom dari tabel kanan otomatis terisi nilai kosong (\texttt{NaN}). Karena kedua DataFrame menggunakan label kunci pengenal yang berbeda (\texttt{Emp\_id} pada tabel kiri dan \texttt{ID} pada tabel kanan), didefinisikan parameter \texttt{left\_on='Emp\_id'} dan \texttt{right\_on='ID'} untuk mencocokkan baris data secara akurat. Selain itu, kolisi atau konflik nama kolom diselesaikan menggunakan parameter \texttt{suffixes=('\_karyawan', '\_kepegawaian')} karena kedua tabel awal sama-sama memiliki atribut bernama \texttt{Status}. Penerapan suffix ini berhasil mengubah label masing-masing kolom menjadi \texttt{Status\_karyawan} dan \texttt{Status\_kepegawaian}, menjaga arti kontekstual masing-masing data tanpa saling menimpa.
    \end{enumerate}
\end{enumerate}
\clearpage

%======================================================================
% BAB IV: KESIMPULAN
%======================================================================
\chapter{Kesimpulan}
Berdasarkan serangkaian langkah praktikum dan penyelesaian tugas pada Pertemuan I, dapat ditarik beberapa kesimpulan penting:
\begin{enumerate}
    \item \textbf{Efisiensi NumPy}: Pustaka \textit{NumPy} dengan objek \textit{ndarray} menyediakan kapabilitas komputasi matriks berkinerja tinggi, manipulasi bentuk array melalui \textit{reshaping} yang hemat memori, serta kemampuan \textit{vectorized boolean filtering} yang sangat efisien untuk seleksi data numerik bervolume besar.
    \item \textbf{Fleksibilitas Pandas}: Pustaka \textit{Pandas} mempermudah manipulasi data tabular melalui antarmuka \textit{DataFrame}, memudahkan pengelolaan indeks, pengurutan, pengelompokan data multidimensi melalui \texttt{groupby()} dan \texttt{agg()}, serta restrukturisasi tabel melalui operasi \textit{pivot} dan \textit{melt}.
    \item \textbf{Integritas Penggabungan Data}: Operasi \texttt{merge()} pada \textit{Pandas} menyediakan fleksibilitas relasional yang setara dengan operasi SQL \textit{JOIN}, di mana penggunaan parameter \texttt{left\_on}, \texttt{right\_on}, dan \texttt{suffixes} terbukti andal dalam menyatukan dataset heterogen tanpa kehilangan konteks maupun duplikasi kolom.
\end{enumerate}
\clearpage

%======================================================================
% DAFTAR PUSTAKA
%======================================================================
\addcontentsline{toc}{chapter}{Daftar Pustaka}
\renewcommand{\bibname}{Daftar Pustaka}
\begin{thebibliography}{99}
    \bibitem{ref1} C. R. Harris, K. J. Millman, S. J. van der Walt, R. Gommers, P. Virtanen, D. Cournapeau, E. Wieser, J. Taylor, S. Berg, N. J. Smith, \textit{et al.}, ``Array programming with NumPy,'' \textit{Nature}, vol. 585, no. 7825, pp. 357--362, 2020.
    \bibitem{ref2} W. McKinney, \textit{Python for Data Analysis: Data Wrangling with pandas, NumPy, and Jupyter}, 3rd ed., Sebastopol, CA: O'Reilly Media, 2022.
    \bibitem{numpy_doc} NumPy Developers, ``NumPy user guide,'' \textit{NumPy v2.0 Manual}, 2024. [Online]. Available: \url{https://numpy.org/doc/}.
    \bibitem{pandas_doc} Pandas Development Team, ``pandas documentation,'' \textit{pandas: powerful Python data analysis toolkit}, 2024. [Online]. Available: \url{https://pandas.pydata.org/docs/}.
\end{thebibliography}

\end{document}
